\scanpage{12}
\mspage{2}

remains covering.

\uline{Hypothesis 2}\unsure{a|.} For $\alpha,\beta\in\mathrm{Ob}\,A$, the set of \deleted{arrows} diagrams
\diagram{\begin{tikzcd}[ampersand replacement=\&,column sep=small,row sep=small]\&\alpha\\\gamma\arrow[ur]\arrow[dr]\&\\\&\beta\end{tikzcd}}
in $A$ (i.e. the set of objects of $A_{/\alpha\times\beta}$, where $\alpha\times\beta$ is the Cartesian product in $\widehat A$), indexing the corresponding arrows $U_\gamma\longrightarrow U_\alpha\times U_\beta$, gives a family \unsure{covering|covering} $U_\alpha\times U_\beta$.

[Likewise for a finite \inserted{(nonempty)} sequence $\alpha_0,\alpha_1,\ldots,\alpha_n$ ($n\in\mathbb N^{\deleted{\unsure{*}}}$) of objects of $A$, and the \deleted{diagrams} arrows
\[U_\gamma\longrightarrow U_{\alpha_0}\times U_{\alpha_1}\times\cdots\times U_{\alpha_n}\]
indexed by the objects of $A_{/\alpha_0\times\alpha_1\times\cdots\alpha_n}$.]
\margin{[Hypothesis 2 generalized\unsure{.|:}]}

For $n=0$, we recover Hypothesis 1', for $n=1$, the initial formulation of Hypothesis 2. But does Hypothesis 2 for $n=1$ imply Hypothesis 2 for \deleted{$\unsure{n>1}$ \unsure{arbitrary}} $n\geq1$\unsure{.|?} We see that this is so, by induction, using the fact that the composite family of two covering families
\diagram{\begin{tikzcd}[ampersand replacement=\&,column sep=large]
U_{\gamma'}\arrow[r,"\varphi(v)"]\&U_\gamma\times U_{\alpha_{n+1}}\arrow[r,"\varphi(u)\times\mathrm{id}_{U_{\alpha_{n+1}}}"]\&(U_{\alpha_0}\times\cdots\times U_{\alpha_n})\times U_{\alpha_{n+\unsure{1}}}
\end{tikzcd}}
($v$ ranging over the arrows $\gamma'\longrightarrow\gamma\times\alpha_{n+1}$, and $u$ over the arrows \unsure{.} $\gamma\longrightarrow\alpha_0\times\cdots\times\alpha_n$) is covering. Thus Hypothesis 1' \deleted{\unsure{implies}} \inserted{$+$} Hypothesis 2\unsure{a|.}
\margin{11}
